\begin{abstract}
空间公共物品博弈中，合作成本由个体承担，收益却由周围成员共享，个体学习因此面临即时利益与群体合作的矛盾。邻居的动作和奖励可以提供额外经验，但报告失真也可能使这种参考损害学习。本文围绕邻居信息进入动作价值更新的方式，研究可靠信息下的合作促进，以及不可靠消息下的参考选择与影响控制。

在可靠信息条件下，本文提出邻居感知价值修正方法：完成标准Q-learning时序差分更新后，根据参考邻居的相对综合奖励与动作一致性，调整自身已执行动作的Q值。前期已发表实验表明，该方法在所研究条件下能够促进合作恢复，其作用取决于感知范围、影响强度与奖励组成。

在不可靠信息条件下，本文以局部消息和自身经验构造线性评分器，选择参考邻居，再通过连续门控调节附加更新幅度。共享参数由完整仿真的交叉熵搜索获得，执行时固定，个体Q表仍从头在线学习。性质分析给出报告动作对合作价值差的作用方向、单步修正幅度界和有界奖励下的Q值范围，解释信息控制的作用与数值边界，不将其作为长期合作保证。

冻结参数后的1200次测试覆盖两个协同因子及干净、持续虚假高奖励背叛报告条件，每种方法每个条件使用30个独立环境种子。完整模型在两个受扰条件的全程合作率为0.96774和0.98398；四个条件下均高于已调强度的固定合作规则，优势为4.105至5.255个百分点，八项主要比较校正后对应的四个区间均为正。

结构比较表明，固定合作评分加学习门控已解释大部分提升；进一步学习评分仅在部分条件增加收益，并在一个干净条件下低于固定评分门控。特征移除显示当前参数依赖显式合作语义，也使用自身经验与价值信息。泛化开发揭示了低协同因子、高异常比例和较小感知范围下的限制。因此，本文所得收益属于明确条件下的任务相关信息控制，高合作不等同于普遍的消息真实性识别。
\end{abstract}

\begin{abstract*}
In spatial public goods games, individuals bear the cost of cooperation while sharing its benefits with nearby players. Individual learning thus faces a tension between immediate returns and collective cooperation. Neighbors' actions and rewards can provide additional experience, but distorted reports may also misdirect learning. This thesis studies how neighbor information enters action-value updates, from promoting cooperation with reliable observations to selecting references and controlling their influence under unreliable messages.

For reliable information, the proposed neighbor-aware correction adjusts the value of the action just taken after the standard Q-learning temporal-difference update. The correction depends on the reference neighbor's relative composite reward and action agreement. Previously published experiments show cooperation recovery under the investigated conditions, with outcomes depending on perception range, influence strength, and reward composition.

For unreliable information, a linear scorer uses local messages and the receiver's own experience to select a reference, while a continuous gate scales the additional update. Shared parameters are optimized by cross-entropy search over complete simulations and remain fixed during evaluation; individual Q-tables continue to learn from scratch. Analysis establishes how reported actions affect the cooperation-value gap, bounds the single-step correction, and derives a Q-value range under bounded rewards. These properties describe the update and its numerical limits rather than guarantee long-term cooperation.

The frozen controllers are tested in 1,200 runs covering two synergy factors under clean messages or persistent false high-reward defection reports, with 30 independent environment seeds per method and condition. The full controller achieves whole-trajectory cooperation of 0.96774 and 0.98398 in the two faulty conditions. It exceeds a tuned fixed cooperative rule by 4.105--5.255 percentage points across all four conditions; the four corresponding intervals remain positive after adjustment over eight primary comparisons.

Structural comparisons show that a learned gate with fixed cooperative scoring accounts for most of the improvement. Learning the scorer provides additional gains only in some conditions and performs worse than fixed scoring with a learned gate in one clean condition. Feature removal reveals dependence on explicit cooperative semantics as well as the receiver's own experience and value information. Development experiments identify limits under lower synergy, higher fault fractions, and smaller perception ranges. The results therefore support task-oriented information control within specified conditions, not general identification of truthful messages.
\end{abstract*}
